% ═══════════════════════════════════════════════════════════
% 14_Route_D_contradiction.tex — v2 CORRIGÉE
% Route D : esquisse contradictoire — NON FERMÉE
% Les failles identifiées par InterIA sont documentées.
% ═══════════════════════════════════════════════════════════

\subsection{Route D~: esquisse contradictoire (non ferm\'ee)}

\paragraph{Objectif.}
Montrer que $\kappa(q_k)\ge\lambda>0$ le long de la tour primoriale,
par contradiction avec les r\'esultats~$\Omega$ de la th\'eorie
analytique des nombres.

\paragraph{Statut.} \textbf{Route~D n'est pas ferm\'ee.}
L'argument ci-dessous identifie trois verrous non r\'esolus
(signal\'es~\S D.4). Il est conserv\'e comme esquisse strat\'egique.

\subsubsection{D.1~: L'identit\'e de d\'epart}

L'observable est~:
\[
\kappa(q)^2 = \frac{1}{\varphi(q)}\sum_{\substack{a\le q\\\gcd(a,q)=1}}M(a)^2.
\]

\begin{remark}[Avertissement : $M(a)$ n'est pas p\'eriodique]\label{rem:not_periodic}
La fonction $a\mapsto M(a)$ d\'epend de l'entier $a$, pas du r\'esidu
$a\bmod q$. Elle n'est \textbf{pas} une fonction sur le groupe
$(\Z/q\Z)^\times$. On ne peut donc pas lui appliquer directement
Parseval sur ce groupe. Toute d\'ecomposition spectrale de $\kappa^2$
doit passer par l'identit\'e d'inclusion--exclusion (Route~C),
pas par l'orthogonalit\'e des caract\`eres.
\end{remark}

\subsubsection{D.2~: Le m\'ecanisme souhait\'e}

L'id\'ee de Route~D est la suivante~:
\begin{enumerate}
\item Supposer $\kappa(q_k)\to 0$.
\item En d\'eduire que $\sum_{\gcd(a,q_k)=1}M(a)^2 = o(\varphi(q_k))$.
\item Montrer que cela contredit un r\'esultat $\Omega$ connu
      sur le comportement de $M(x)$.
\end{enumerate}

\subsubsection{D.3~: La contradiction tent\'ee}

\paragraph{R\'esultat $\Omega$ de type Ingham.}
De mani\`ere inconditionnelle (z\'eros critiques de $\zeta$,
Hardy 1914)~:
\[
\sum_{n\le x}\mu(n) = \Omega(x^{1/2}).
\]
Autrement dit, $|M(x)|\ge c\cdot x^{1/2}$ pour une infinit\'e
de $x\to\infty$.

\paragraph{Tentative de transfert.}
Si ces grandes valeurs de $M(x)$ se produisent \`a des $x$
copremiers \`a $q_k$, alors~:
\[
\kappa(q_k)^2 \ge \frac{1}{\varphi(q_k)}\cdot c^2 q_k
= \frac{c^2 q_k}{\varphi(q_k)} \to \infty.
\]
Contradiction avec $\kappa\to 0$.

\subsubsection{D.4~: Les trois verrous non r\'esolus}

\begin{remark}[Verrou 1 : transfert aux copremiers]\label{rem:V1}
La densit\'e des entiers copremiers \`a $q_k$ est
$\varphi(q_k)/q_k \sim e^{-\gamma}/\log p_k \to 0$.
Les grandes valeurs de $M(x)$ pourraient se produire exclusivement
\`a des $x$ non copremiers \`a $q_k$. Le transfert n\'ecessite de
montrer que $M$ atteint $\Omega(x^{1/2})$ sur les copremiers,
ce qui n'est pas garanti par le r\'esultat $\Omega$ standard.
\end{remark}

\begin{remark}[Verrou 2 : force de la minoration]\label{rem:V2}
M\^eme avec un seul $a$ copremier tel que $|M(a)|\ge c\sqrt{q}$,
on obtient seulement~:
\[
\kappa^2 \ge \frac{c^2 q}{\varphi(q)} = O\!\left(\frac{\log\log q}{1}\right),
\]
ce qui prouve $\kappa\to\infty$ (si le transfert fonctionne),
pas $\kappa\ge\lambda>0$.

\textbf{Attention}~: si on n'obtient qu'un r\'esultat plus faible
$|M(a)|=\Omega(q^{1/2-\epsilon})$ sur les copremiers, la minoration
donne seulement $\kappa^2\ge\Omega(q^{-2\epsilon}\log\log q)$,
qui peut encore tendre vers z\'ero.
\end{remark}

\begin{remark}[Verrou 3 : $M$ vs $\mu$]\label{rem:V3}
Les outils analytiques (Perron, formule explicite) s'appliquent
naturellement \`a $\sum\mu(n)\chi(n)$ (li\'e \`a $1/L(s,\chi)$),
pas \`a $\sum M(n)\chi(n)$. Le passage de $\mu$ \`a $M$ par
sommation d'Abel introduit un terme de convolution qui peut
d\'etruire les estimations ponctuelles.
\end{remark}

\subsubsection{D.5~: Ce qui serait suffisant}

Pour fermer Route~D, il suffirait de prouver l'un des \'enonc\'es
suivants~:

\begin{enumerate}
\item \textbf{(Transfert fort)} Pour tout primorial $q$ assez grand,
il existe $a\le q$ avec $\gcd(a,q)=1$ et $|M(a)|\ge c\sqrt{q}$.

\item \textbf{(Second moment copremier)} $\sum_{\gcd(a,q)=1}M(a)^2
\ge c\cdot q^2/\varphi(q)$ pour une constante $c>0$.
\emph{Note}~: c'est exactement Route~C raffin\'ee.

\item \textbf{(Rubinstein--Sarnak copremier)} Sous RH+LI,
la distribution de $M(a)/\sqrt{a}$ restreinte aux copremiers
a une variance born\'ee inf\'erieurement.
\end{enumerate}

\subsubsection{D.6~: Tableau comparatif mis \`a jour}

\begin{center}
\begin{tabular}{lllll}
\toprule
\textbf{Route} & \textbf{\'Enonc\'e} & \textbf{Hypoth\`ese} & \textbf{Force} & \textbf{Statut} \\
\midrule
C na\"ive & $S_1$ domine & Titchmarsh & $\kappa\to\infty$ & R\'efut\'ee \\
C raffin\'ee & $\sum\mu(d)S_d>0$ & Convexit\'e & $\kappa>0$ & Ouverte \\
D faible & $\kappa\not\to 0$ & Transfert & $\liminf>0$ & \textbf{Ouverte (3 verrous)} \\
D+RS & $\kappa\ge\lambda$ & RH+LI & $\inf>0$ & Conditionnelle \\
A & bloc local & MR 2016 & $\kappa\ge\lambda$ & Ouverte \\
\bottomrule
\end{tabular}
\end{center}

\begin{remark}[Fait empirique]
Malgr\'e l'absence de preuve, les donn\'ees num\'eriques sont
sans ambigu\"it\'e~: $\kappa\sim q^{0.337}$, $R^2=0.96$,
strictement croissant sur 7 niveaux primoriaux.
Le m\'ecanisme analytique sous-jacent reste \`a identifier.
\end{remark}
